% Generated draft: requires mathematical and source review.
\documentclass[11pt]{article}
\usepackage[margin=25mm]{geometry}
\usepackage{fontspec}
\setmainfont{Latin Modern Roman}
\usepackage{amsmath,amssymb,mathtools,unicode-math}
\setmathfont{Latin Modern Math}
\usepackage{enumitem,needspace,xurl,hyperref}
\setlength{\parindent}{0pt}
\setlength{\parskip}{5pt}
\setlength{\emergencystretch}{3em}
\newcommand{\N}{\mathbb{N}}
\newcommand{\Z}{\mathbb{Z}}
\newcommand{\Q}{\mathbb{Q}}
\newcommand{\R}{\mathbb{R}}
\newcommand{\C}{\mathbb{C}}
\newcommand{\F}{\mathbb{F}}
\newcommand{\A}{\mathbb{A}}
\newcommand{\PP}{\mathbb{P}}
\newcommand{\E}{\mathbb{E}}
\newcommand{\eps}{\varepsilon}
\newcommand{\dd}{\,\mathrm{d}}
\begin{document}
\section*{UnsolvedMath ID: 800636}
\textbf{Source problem code:} catalogue entry
\subsection{Definitions and mathematical statement}
Let \(T\) be an ambient topological space and give each \(X\subseteq T\) its subspace topology. The ambient space, permitted subspaces, and cover conventions must be specified to obtain a definite formulation; the supplied record does not specify them.

Under the usual open-cover interpretation, an admissible cover \(\mathcal U\) of \(X\) consists of open subsets of \(X\), satisfies \(\bigcup\mathcal U=X\), and does not contain \(X\) itself. An \(\omega\)-cover is such a cover satisfying
\[
\forall F\subseteq X\;\bigl(F\text{ finite}\Longrightarrow \exists U\in\mathcal U\;(F\subseteq U)\bigr).
\]
A \(\gamma\)-cover is an infinite such cover satisfying
\[
\forall x\in X,\qquad \{U\in\mathcal U:x\notin U\}\text{ is finite}.
\]
Write \(\Omega(X)\) and \(\Gamma(X)\) for the corresponding families of admissible covers. Whether admissible covers must be countable is unresolved in the supplied record; these definitions are to be restricted accordingly once that convention is established.

For families of covers \(\mathcal A(X)\) and \(\mathcal B(X)\), the property \(\mathsf S_1(\mathcal A,\mathcal B)\) means that for every sequence \((\mathcal U_n)_{n\geq1}\) from \(\mathcal A(X)\), there are choices \(U_n\in\mathcal U_n\) such that \(\{U_n:n\geq1\}\in\mathcal B(X)\). The property \(\mathsf S_{\mathrm{fin}}(\mathcal A,\mathcal B)\) means that for every such sequence there are finite subfamilies \(\mathcal F_n\subseteq\mathcal U_n\) such that
\[
\bigcup_{n\geq1}\mathcal F_n\in\mathcal B(X).
\]
Here the displayed union is a family of open sets.

\textbf{Question (source Problem 8, \texttt{[AddQuad]}).} Is at least one of the three properties
\[
\mathsf S_{\mathrm{fin}}(\Omega,\Omega),\qquad
\mathsf S_1(\Gamma,\Omega),\qquad
\mathsf S_{\mathrm{fin}}(\Gamma,\Omega)
\]
consistently closed under finite unions? For each candidate property \(P\), the finite-union assertion, once the intended ambient setting is fixed, is
\[
\forall m\geq1\;\forall X_1,\ldots,X_m\subseteq T,\qquad
\left(\bigwedge_{j=1}^{m}P(X_j)\right)
\Longrightarrow
P\!\left(\bigcup_{j=1}^{m}X_j\right),
\]
with the quantifiers restricted to the intended permitted subspaces. Under the usual set-theoretic reading, the question asks whether this assertion for at least one candidate is consistent with \(\mathrm{ZFC}\). Neither that base theory nor the intended ambient setting is explicit in the supplied record. Thus this is a formulation schema with unresolved qualifications, not a uniquely specified consistency assertion.
\subsection{Short English statement}
The source asks whether any of three selection properties can consistently be preserved when finitely many spaces having that property are united. A precise interpretation still requires the intended class of spaces, cover countability conventions, and consistency base theory. The registry labels the problem open; this draft does not independently verify that status.
\subsection{Sources}
\begin{itemize}
\item Local UnsolvedMath registry, permanent ID \(800636\), exact source problem code \texttt{catalogue entry}; supplied title: \emph{Selection Principles Milestone Problems --- Problem 8 ([AddQuad])}. The record attributes the problem to Tsaban, with proposed year \(2003\), and supplies status \texttt{open}.
\item Tsaban, \emph{Selection Principles in Mathematics: A Milestone of Open Problems} (2003), Problem 8. Supplied source URL: \url{https://arxiv.org/abs/math/0312182}. The record reports extraction by \texttt{arxiv\_latex} and access on \texttt{2026-10-16}; these are supplied provenance details, not an independently verified visit or literature assessment.
\item The supplied attribution/context marker is \texttt{[AddQuad]}. Its bibliographic expansion is not included in the record.
\end{itemize}
\end{document}
